\documentclass[12pt, a4paper]{article}

% --- Packages ---
\usepackage[margin=1in]{geometry}
\usepackage{amsmath}
\usepackage{hyperref}
\usepackage{enumitem}
\usepackage{verbatim}
\usepackage{graphicx}
\usepackage{xcolor}
\usepackage{mathptmx}
\usepackage[T1]{fontenc}
\usepackage{setspace}
\usepackage{needspace}
\usepackage{titlesec}
\usepackage{etoolbox}

% --- Typographical Adjustments to Prevent Awkward Page Breaks ---
\raggedbottom
\widowpenalty=10000
\clubpenalty=10000
\interlinepenalty=100
\setstretch{1.15}

% --- Section heading formatting ---
\titleformat{\section}{\normalfont\Large\bfseries}{\thesection}{1em}{}
\titlespacing*{\section}{0pt}{2ex plus 1ex minus .2ex}{1ex plus .2ex}

\titleformat{\subsection}{\normalfont\large\bfseries}{\thesubsection}{1em}{}
\titlespacing*{\subsection}{0pt}{1.5ex plus 0.5ex minus .2ex}{0.75ex plus .2ex}

% Safe needspace hook that works for both \section and \section* ---
\makeatletter
\pretocmd{\@startsection}{\needspace{5\baselineskip}}{}{}
\makeatother

% Suppress the auto-heading thebibliography generates; we place our own below
\renewcommand{\refname}{}

% ============================================================
\begin{document}

% ==================== TITLE PAGE ====================
\begin{titlepage}
    \centering
    \vspace*{1cm}

    % Logo (ensure giki-logo-1-e1659884838448 image file is in the same folder as this .tex file)
    \includegraphics[width=0.25\textwidth]{giki-logo-1-e1659884838448}\\[1.5cm]

    {\color{blue}\bfseries\LARGE Algorithm Design and Implementation for\\[0.3cm]
    Automated Timetable Scheduling}\\[0.6cm]

    {\Large \textbf{Semester Project Report}}\\[0.4cm]
    {\Large \textbf{CS378: Design and Analysis of Algorithms}}\\[0.3cm]
    {\Large Spring 2026}\\[1.5cm]

    {\Large \textbf{Instructor: Salman Ashraf}}\\[1.8cm]

    {\Large \textbf{Submitted By:}}\\[0.4cm]

    \begin{tabular}{ll}
        Aazeb Ali      & 2023003 \\
        Abdullah Ihsan & 2023039 \\
        Rabbin Batool  & 2023586 \\
        Zain Jamshaid  & 2023772 \\
    \end{tabular}\\[2.5cm]

    {\large Faculty of Computer Science and Engineering}\\[0.3cm]
    {\large GIK Institute of Engineering Sciences and Technology}\\[0.5cm]

    \vfill
    {\large \textbf{Semester: Spring 2026}}
    \vspace{1cm}
\end{titlepage}
% ====================================================

\newpage

% --- Abstract ---
\begin{abstract}
\noindent The University Course Timetabling Problem (UCTP) is a well-known NP-Hard
optimization challenge that requires mapping a set of academic events to specific times
and spatial resources without violating complex structural constraints. This report details
the design, mathematical analysis, and implementation strategy of a hybrid algorithm
utilizing Greedy Construction combined with Backtracking Swap Repair. The proposed
solution is specifically tailored to the operational requirements of the Ghulam Ishaq Khan
Institute (GIKI), balancing strict administrative mandates (hard constraints) with
pedagogical optimization goals (soft constraints) to generate a robust, conflict-free
weekly schedule.
\end{abstract}

\vspace{0.5cm}
\hrule
\vspace{0.5cm}

% --- Section 1: Problem Definition ---
\section{Problem Definition}

The timetable scheduling problem involves assigning a defined set of academic
entities---specifically courses, instructors, student cohorts, and classrooms---into a
finite grid of weekly time slots. The objective is to produce a feasible schedule where
no entities overlap in an invalid manner, and resource utilization is maximized.

Formally, the system must process:

\begin{itemize}[noitemsep, topsep=4pt]
    \item \textbf{Entities:} A set of $N$ course sections ($S$) and a set of $M$
          available classrooms ($R$). Each section requires a specific number of
          instructional hours, and each room possesses a maximum seating capacity.

    \item \textbf{Time Framework:} A standard 5-day academic week partitioned into
          distinct daily slots, incorporating a mandatory lunch break.

    \item \textbf{Objective:} To map $S \rightarrow (R \times T)$, where $T$ represents
          the set of available time slots, such that all hard constraints are strictly
          satisfied, and soft constraints are optimally managed to ensure workload
          balance.
\end{itemize}

% --- Section 2: Algorithm Design ---
\section{Algorithm Design}

Given that standard brute-force methodologies possess an exponential time complexity
($O(k^N)$), they are strictly computationally infeasible for a university-scale dataset.
Therefore, this project implements a hybrid heuristic algorithm: \textbf{Greedy
Construction combined with Depth-1 Backtracking Swap Repair and Progressive
Relaxation}.

The algorithm operates in a multi-phased pipeline:

\begin{itemize}[topsep=4pt]
    \item \textbf{Phase 0: Heuristic Sorting (Constraint Density):} Prior to assignment,
          all sections are sorted by ``constraint density'' using four ordered keys:
          (1)~named instructor before \texttt{TBA} (real scheduling constraints take
          precedence); (2)~non-elective before elective (core courses are harder to
          place and must not be crowded out); (3)~higher credit hours first; and
          (4)~more programme codes first (a section shared across multiple programmes
          is the most constrained). This ``hardest-first'' ordering ensures that
          flexible, low-constraint sections fill the remaining gaps rather than
          monopolising critical resources.

    \item \textbf{Phase 1: Strict Greedy Construction:} The algorithm iterates
          sequentially through the sorted sections, attempting to place each course into
          the first valid slot-room combination that satisfies all constraints. A tracking
          array (\texttt{dayLoad}) is maintained to ensure classes are balanced across
          the week.

    \item \textbf{Phase 2: Backtracking Swap Repair:} For each section $U$ that
          fails Phase~1, the algorithm first re-attempts direct greedy placement
          (previously scheduled sections may have opened new slot combinations).
          If that still fails, a Depth-1 Local Search is triggered: a candidate
          victim $V$ that conflicts with $U$ (same instructor, or overlapping
          programme and year level) is temporarily removed from the schedule;
          $U$ is placed in the freed slot; and the algorithm attempts to re-home $V$
          elsewhere. If $V$ cannot be re-placed, the swap is fully undone and the
          next candidate victim is tried. This resolves localised deadlocks without
          spiraling into infinite recursive loops.

    \item \textbf{Phases 3--6: Progressive Constraint Relaxation:} For sections that
          remain unscheduled, the algorithm walks through four escalating relaxation
          steps. \textbf{Phase~3} relaxes H3 for elective courses, permitting them
          to overlap with other programmes. \textbf{Phase~4} further relaxes H3 for
          distinct section groups within the same programme (e.g.\ Section~A and
          Section~B of the same course may coexist in the same slot). \textbf{Phase~5}
          skips H2 entirely and appends ``\texttt{(TBD)}'' to the instructor name,
          flagging the assignment for manual instructor resolution. \textbf{Phase~6}
          skips both H2 and H3, tagging the assignment as ``\texttt{(FORCE)}'' as an
          absolute last resort. Together these phases guarantee a complete schedule is
          always produced.
\end{itemize}

% --- Section 3: Assumptions ---
\section{Assumptions}

To mathematically model the GIKI academic environment, the following operational
assumptions were formulated:

\begin{enumerate}[label=\arabic*., noitemsep, topsep=4pt]
    \item The academic week consists of 5 working days (Monday to Friday).

    \item There are 8 operational time slots per day, starting at 08:30.
          The daily schedule runs from 08:30 to 17:20, partitioned as:
          \textbf{Morning} (slots 0--4): 08:30, 09:30, 10:30, 11:30, 12:30;
          \textbf{Afternoon} (slots 5--7): 14:30, 15:30, 16:30.

    \item A mandatory break exists between 13:30 and 14:30 (Prayer /
          Jumu'ah break). The last morning slot (slot~4) ends at 13:20 and
          the first afternoon slot (slot~5) starts at 14:30, so no class
          is ever scheduled during the break window.
          Laboratory sessions must be placed entirely within the morning
          block or entirely within the afternoon block.

    \item Standard 3-credit hour (CH) courses are mapped to
          $3\ \text{days} \times 1\ \text{hour}$, preferring balanced spreads
          (e.g., MWF).

    \item Standard 2-CH courses are mapped to
          $2\ \text{days} \times 1\ \text{hour}$, preferring TTh configurations.

    \item A 1-CH laboratory session requires 3 consecutive hours on a single day and
          cannot straddle the lunch break.

    \item Default course enrollment is assumed to be 40.

    \item Specific administrative venues (e.g., ES LH4) are reserved for centralized
          assessments and excluded from routine scheduling.

    \item Parallel course sections (e.g., Section A vs.\ Section B) serve distinct
          student cohorts and may overlap chronologically.

    \item Back-to-back teaching is permitted. Hard constraint H2 already
          prevents an instructor from being double-booked in overlapping
          slots; no additional minimum-gap restriction is imposed. This
          ensures that all 8 daily time slots remain reachable for every
          instructor, avoiding artificially empty rows in the schedule.
\end{enumerate}

% --- Section 4: Constraint Handling Strategy ---
\section{Constraint Handling Strategy}

\subsection{Hard Constraints (Strict Compliance)}

\begin{itemize}[noitemsep, topsep=4pt]
    \item \textbf{Room Uniqueness (H1):} Addressed by spatial checking loops. A
          classroom cannot host multiple sections in the same time slot.

    \item \textbf{Instructor Uniqueness (H2):} Evaluated by validating the instructor's
          timeline. An instructor cannot be assigned to overlapping temporal slots.

    \item \textbf{Program/Cohort Conflicts (H3):} Solved via set-intersection logic
          (\texttt{programsOverlap}). Sections belonging to the same academic program
          and year level cannot be scheduled simultaneously, preventing cohort overlaps.

    \item \textbf{Room Capacity (H4):} Enforced during the initial sorting logic
          (\texttt{sortRooms}). The algorithm filters out any room where capacity is
          strictly less than the section's enrollment.
\end{itemize}

\subsection{Soft Constraints (Optimization Goals)}

\begin{itemize}[noitemsep, topsep=4pt]
    \item \textbf{Workload Distribution:} Managed via the \texttt{dayLoad} array. The
          algorithm actively monitors daily class distributions and prioritizes
          combinations that place courses on the least congested days, ensuring
          no single weekday becomes disproportionately overloaded.

    \item \textbf{Efficient Room Allocation:} Handled via a two-key heuristic
          (\texttt{sortRooms}). The primary key is \textit{departmental ownership}:
          rooms whose allocation field matches the section's programme are listed
          first. The secondary key is \textit{best-fit capacity}: among eligible rooms,
          the one with capacity closest to (but $\geq$) the section's enrollment is
          selected, minimising wasted seating and preventing large halls from being
          monopolised by small cohorts.
\end{itemize}

% --- Section 5: Input and Output Specification ---
\section{Input and Output Specification}

\noindent\textbf{Input Data:}\\[4pt]
The system ingests structured data representing two primary sets:

\begin{itemize}[noitemsep, topsep=4pt]
    \item $S$ (Sections Data): Includes \texttt{courseCode}, \texttt{courseTitle},
          \texttt{sectionLabel}, \texttt{program}, \texttt{instructorName},
          \texttt{creditHours}, \texttt{isLab}, \texttt{isElective},
          \texttt{sectionGroup}, \texttt{yearLevel}, and \texttt{enrollment}.
          The \texttt{sectionGroup} identifier (e.g.\ ``A'', ``B'') allows Phase~4
          to distinguish parallel sections serving distinct cohorts.
          \texttt{yearLevel} (1--4) is not present in the Excel files; it is
          \textit{automatically inferred} from the course code during seeding
          (e.g.\ \texttt{CS232} $\Rightarrow$ year~2, \texttt{CV314} $\Rightarrow$
          year~3) via the \texttt{extractYearLevel()} helper, which reads the first
          digit of the numeric part of the code. A value of~0 is stored when no
          valid digit is found.

    \item $R$ (Rooms Data): Includes \texttt{name}, \texttt{capacity}, and
          departmental \texttt{allocation}.
\end{itemize}

\noindent\textbf{Output Data:}\\[4pt]
The primary output is a graphical timetable mapping, alongside an internal data
structure containing:

\begin{itemize}[noitemsep, topsep=4pt]
    \item \texttt{assignments}: A complete list of all mapped sections, dictating the
          assigned \texttt{room} and specific temporal \texttt{slots}.

    \item \texttt{unassigned}: A list of sections the algorithm failed to place entirely
          (flagged for manual review).

    \item \texttt{hardViolations}: An integer metric indicating the overall ``health''
          and validity of the generated schedule (should ideally equal 0).
\end{itemize}

% --- Section 6: Complexity Analysis ---
\section{Complexity Analysis}

\subsection{Time Complexity}

Let $N$ denote the total number of course sections, $M$ denote the total number of
rooms, $C$ denote the maximum possible time combinations for a given section, and $U$
denote the number of unassigned sections pushed to Phase~2.

$C$ is bounded by the 3-CH and 2-CH course types, both of which produce the same
maximum:
$$C = \binom{5}{3} \times 8 = 10 \times 8 = 80 \quad \text{(3-CH: 3 days from 5, any of 8 slots)}$$
$$C = \binom{5}{2} \times 8 = 10 \times 8 = 80 \quad \text{(2-CH: 2 days from 5, any of 8 slots)}$$
Laboratory sessions generate only $4 \times 5 = 20$ combinations (4 valid start slots
$\times$ 5 days), so $C \approx 80$ is the worst-case upper bound across all course types.

\begin{itemize}[topsep=4pt]
    \item \textbf{Phase 1 (Greedy Construction):} For each of the $N$ sections, the
          algorithm tests up to $C$ combinations across $M$ rooms. Every test validates
          conflicts against a maximum of $N$ already-placed sections. Complexity:
          $O(N \cdot C \cdot M \cdot N) = O(N^2 \cdot C \cdot M)$.

    \item \textbf{Phase 2 (Backtracking Repair):} Iterates over $U$ failures. For each,
          it tests swapping against up to $N$ existing assignments, effectively running
          the placement logic twice per swap. Complexity: $O(U \cdot N^2 \cdot C \cdot M)$.
\end{itemize}

Since $U \ll N$ (typically under 5\% of sections fall to Phase 2), $U$ and $C$ act as
constants.

\medskip
\noindent\textbf{Overall Time Complexity:} $\mathbf{O(N^2 \cdot M)}$.

\medskip
\noindent This polynomial boundary confirms the algorithm is highly scalable and will
execute in mere seconds on modern hardware, satisfying the requirement for efficiency.

\subsection{Space Complexity}

The algorithm requires memory proportional to the size of the inputs and the output
structures.

\begin{itemize}[noitemsep, topsep=4pt]
    \item Storing sections $S$ and rooms $R$ requires $O(N + M)$ space.

    \item The \texttt{assignments} array scales linearly with the number of placed
          sections: $O(N)$.

    \item Due to restricting the backtracking to a ``Depth-1'' limit, the call stack
          does not grow recursively, ensuring minimal overhead memory usage.
\end{itemize}

\medskip
\noindent\textbf{Overall Space Complexity:} $\mathbf{O(N + M)}$.

% --- Section 7: Limitations ---
\section{Limitations}

While highly effective, the current design exhibits specific limitations:

\begin{enumerate}[label=\arabic*., topsep=4pt]
    \item \textbf{The ``Domino Effect'' Restriction:} Because the Local Search is
          strictly Depth-1, it cannot resolve complex deadlocks that require cascading
          moves (e.g., moving Course~A to move Course~B to move Course~C).

    \item \textbf{Global Load Drift:} The \texttt{dayLoad} array globally balances the
          week, but it does not track spatial congestion (e.g., balancing the load
          specifically for one faculty building versus another).

    \item \textbf{Optimal vs.\ Feasible:} As a heuristic approach, the algorithm
          guarantees a \textit{feasible} (good enough) schedule rapidly, but it cannot
          mathematically guarantee the absolute \textit{optimal} (perfect) global
          minimum solution.
\end{enumerate}

% --- Section 8: Results and Sample Outputs ---
\section{Results and Sample Outputs}

\textit{(Note: Screenshots of the implementation should be inserted here using the
\texttt{\textbackslash includegraphics} command.)}

\vspace{0.3cm}
\noindent The algorithmic core was connected to a Graphical User Interface (GUI). The
interface successfully demonstrates:

\begin{itemize}[topsep=4pt]
    \item \textbf{Data Parsing Screen:} Users can upload or define courses, capacities,
          and teachers.

    \item \textbf{Execution Metrics:} Upon running the algorithm, the GUI displays the
          execution time (in milliseconds) and the total count of Hard Violations.

    \item \textbf{Visual Grid:} The primary output is a color-coded weekly calendar
          view. The X-axis represents the 5 working days, and the Y-axis represents the
          8 daily time slots. Users can filter this grid by Room, Instructor, or Student
          Cohort to verify constraints visually.
\end{itemize}

% --- Section 9: Conclusion ---
\section{Conclusion}

The implemented algorithm successfully bridges the gap between strict theoretical
constraints and real-world administrative compromises at GIK Institute. By leveraging a
Greedy framework fortified with Constraint Density Sorting and Depth-1 Backtracking,
the system achieves a computationally efficient ($O(N^2 \cdot M)$) solution to the
NP-Hard Timetabling Problem. The strategic relaxation of constraints in the final phases
ensures a complete schedule is generated, demonstrating practical, software-engineering-level
algorithm design capable of automating complex institutional logistics.

% --- Project Links ---
\section{Project Links}

\begin{itemize}[noitemsep, topsep=4pt]
    \item \textbf{Live Demo:}
          \href{https://algo-project-five.vercel.app/}{\texttt{https://algo-project-five.vercel.app/}}

    \item \textbf{GitHub Repository:}
          \href{https://github.com/ZainJ5/algo\_project}{\texttt{https://github.com/ZainJ5/algo\_project}}
\end{itemize}

% --- References ---
\newpage
\section{References}

\begin{thebibliography}{9}

\bibitem{cormen}
Cormen, T.\,H., Leiserson, C.\,E., Rivest, R.\,L., \& Stein, C. (2009).
\textit{Introduction to Algorithms} (3rd ed.). MIT Press.

\bibitem{schaerf}
Schaerf, A. (1999).
\textit{A survey of automated timetabling}.
Artificial Intelligence Review, 13(2), 87--127.

\bibitem{carter}
Carter, M.\,W., \& Laporte, G. (1998).
\textit{Recent developments in practical course timetabling}.
Practice and Theory of Automated Timetabling II. Springer.

\end{thebibliography}

\newpage

% --- Appendix: Algorithm / Pseudocode ---
\appendix
\section{Appendix: Step-by-Step Algorithm (Class Format)}

The following pseudocode demonstrates the exact step-by-step logic of the hybrid
scheduling algorithm, structured according to the standard algorithmic notation
discussed in class.

\begin{verbatim}
Algorithm TimetableScheduling(S, n, R, m)
{
    // S: Array of n course sections
    // R: Array of m available rooms

    // STEP 0: Initialization and Sorting
    Assignments = {};
    Unassigned  = {};
    dayLoad[5]  = {0, 0, 0, 0, 0};

    SortByConstraintDensity(S, n);
    // Sorts S so largest enrollments and strict constraints come first

    // STEP 1: Phase 1 - Strict Greedy Construction
    for (i = 0; i < n; i++)
    {
        placed = false;

        // Fetch time combinations (credit hours, lab flag, current day-load)
        combos = GetSlotCombinations(S[i].credits, S[i].isLab, dayLoad);

        for (c = 0; c < combos.length; c++)
        {
            for (j = 0; j < m; j++)
            {
                // Check Hard Constraints (H1, H2, H3, H4)
                if (CheckConflicts(S[i], R[j], combos[c], Assignments) == 0)
                {
                    Assignments.add(S[i], R[j], combos[c]);
                    Update(dayLoad);
                    placed = true;
                    break;
                }
            }
            if (placed == true) { break; }
        }

        if (placed == false) { Unassigned.add(S[i]); }
    }

    // STEP 2: Phase 2 - Backtracking Swap Repair (Depth-1)
    for (u = 0; u < Unassigned.length; u++)
    {
        // First: re-attempt direct placement (schedule may have shifted)
        if (TryPlace(Unassigned[u], R, Assignments, dayLoad) == true)
        {
            Assignments.add(Unassigned[u]);
            Update(dayLoad);
            continue;
        }

        for (v = 0; v < Assignments.length; v++)
        {
            // Only attempt swap if removing V logically helps U
            if (WouldBenefitSwap(Unassigned[u], Assignments[v]) == true)
            {
                TemporarilyRemove(Assignments[v]);

                if (TryPlace(Unassigned[u], R, Assignments) == true)
                {
                    if (TryPlace(Assignments[v], R, Assignments) == true)
                    {
                        break; // Swap successful, move to next unassigned
                    }
                    else
                    {
                        UndoPlacement(Unassigned[u]);
                        RestorePlacement(Assignments[v]);
                    }
                }
                else
                {
                    RestorePlacement(Assignments[v]);
                }
            }
        }
    }

    // STEP 3: Phases 3 to 6 - Progressive Constraint Relaxation
    for (u = 0; u < Unassigned.length; u++)
    {
        // Phase 3: Relax H3 for electives (allow elective program overlaps)
        if (TryPlaceRelaxed(Unassigned[u], relaxElectives=true) == true)
            { continue; }

        // Phase 4: Relax H3 for section groups (Section A vs Section B)
        if (TryPlaceRelaxed(Unassigned[u], relaxElectives=true,
                             relaxGroups=true) == true)
            { continue; }

        // Phase 5: Skip H2, mark instructor as "(TBD)" for manual review
        if (TryPlaceRelaxed(Unassigned[u], skipH2=true,
                             maskInstructor="(TBD)") == true)
            { continue; }

        // Phase 6: Last resort — skip H2+H3, tag instructor as "(FORCE)"
        ForcePlace(Unassigned[u], Assignments);
    }

    return Assignments;
}
\end{verbatim}

\end{document}